\section{Physics-Guided TTA Variants}
\label{sec:method}

This section defines the three physics-guided variants evaluated in the
benchmark. They are not the paper's main contribution; they exist to test
whether a coarse OFDM residual prior can make TTA less harmful than generic
entropy minimisation. The objective is \emph{feature-statistic alignment using a
coarse OFDM residual-prior representation}, not a direct PDE residual loss and
not a channel-estimation method.

\paragraph{Residual-prior representation.}
For an input CSI observation $\mathbf{x}$, we compute an analytical prior
$\mathbf{p}(\mathbf{x})$ from the synthetic OFDM propagation model and feed the
encoder the concatenation
$\tilde{\mathbf{x}}=[\mathbf{x},\mathbf{p}(\mathbf{x}),\mathbf{x}-\mathbf{p}(\mathbf{x})]$.
The prior is intentionally coarse on real data because the true path gains,
delays, and room geometry are unobserved. The classifier head is frozen during
adaptation; only the encoder-side parameters listed in
Appendix~\ref{app:hyperparams} are updated.

\paragraph{Feature-statistic alignment objective.}
After source training, we store source feature mean/variance
$(\boldsymbol{\mu}_s,\boldsymbol{\sigma}_s^2)$ on the encoder representation
$\mathbf{z}=f_\phi(\tilde{\mathbf{x}})$. In the implementation the same
$\mathbf{z}$ is used for both the primary and auxiliary moment terms
(\texttt{outputs["features"]}); we reserve the word ``residual-branch'' for
the input concatenation and do not claim a separate internal residual stream.
We also store auxiliary
statistics $(\boldsymbol{\mu}^{r}_s,(\boldsymbol{\sigma}^{r}_s)^2)$. For an
unlabelled target batch $B_t$ the objective is
\[
    \mathcal{L}_{\text{align}}
    = \|\boldsymbol{\mu}_t-\boldsymbol{\mu}_s\|_2^2
    + \|\boldsymbol{\sigma}_t^2-\boldsymbol{\sigma}_s^2\|_2^2
    + \lambda_r\!\left(
        \|\boldsymbol{\mu}^{r}_t-\boldsymbol{\mu}^{r}_s\|_2^2
        + \|(\boldsymbol{\sigma}^{r}_t)^2-(\boldsymbol{\sigma}^{r}_s)^2\|_2^2
      \right),
\]
with $\lambda_r=1$. \texttt{physics\_tta} applies this objective under the
standard budget. \texttt{safe\_physics} rejects a step when the unlabelled
objective increases or feature-statistic drift exceeds $\tau_a=2.0$.
\texttt{selective\_physics} softly blends the adapted and source predictions,
$q(y|\mathbf{x})=g(\mathbf{x})q_{\text{adapt}}+(1-g(\mathbf{x}))q_{\text{src}}$,
with
$g(\mathbf{x})=\mathrm{clip}((c_{\text{adapt}}-c_{\text{src}})/\tau_c+0.5,0,1)\exp(-d_{\text{align}}/\tau_a)$
where $c$ is max softmax confidence and $d_{\text{align}}$ is feature-statistic
divergence from the source reference. Thresholds ($\tau_c{=}0.2$,
$\tau_a{=}2.0$) were fixed before evaluation.
